% Generated draft: requires mathematical and source review.
\documentclass[11pt]{article}
\usepackage[margin=25mm]{geometry}
\usepackage{fontspec}
\setmainfont{Latin Modern Roman}
\usepackage{amsmath,amssymb,mathtools,unicode-math}
\setmathfont{Latin Modern Math}
\usepackage{enumitem,needspace,xurl,hyperref}
\setlength{\parindent}{0pt}
\setlength{\parskip}{5pt}
\setlength{\emergencystretch}{3em}
\newcommand{\N}{\mathbb{N}}
\newcommand{\Z}{\mathbb{Z}}
\newcommand{\Q}{\mathbb{Q}}
\newcommand{\R}{\mathbb{R}}
\newcommand{\C}{\mathbb{C}}
\newcommand{\F}{\mathbb{F}}
\newcommand{\A}{\mathbb{A}}
\newcommand{\PP}{\mathbb{P}}
\newcommand{\E}{\mathbb{E}}
\newcommand{\eps}{\varepsilon}
\newcommand{\dd}{\,\mathrm{d}}
\begin{document}
\section*{UnsolvedMath ID: 20000631}
\textbf{Source problem code:} catalogue entry
\subsection{Definitions and mathematical statement}
Bernard Shiffman's source question asks whether the Fubini--Study metric on \(\PP^1\) minimizes the expected number of critical points, and proposes an affirmative answer. It gives no degree, probability law, admissible metric class, or definition of critical point. The following is the Douglas--Shiffman--Zelditch-type reconstruction used by the attached aid, not a verified complete statement of the original question.

Let \(h\) be a smooth positive Hermitian metric on \(\mathcal O(1)\to\PP^1\), with positive curvature form \(\omega_h\) in a fixed area class. For an integer \(N\geq1\), equip \(H^0(\PP^1,\mathcal O(N))\) with the \(L^2\) inner product from \(h^N\) and its curvature volume form. In an orthonormal basis choose independent standard complex Gaussian coefficients for a random section \(s\). A Chern-critical point is a zero of \(\nabla^{1,0}_{h^N}s\); for sections with simple zeros these are the nonzero critical points of the pointwise norm, so the zeros of \(s\) themselves are not added to the count. Define
\[
\mathcal N_N(h)=\mathbb E_h\#\operatorname{Crit}_{h^N}(s).
\]
The reconstructed fixed-degree minimization asks whether
\[
\mathcal N_N(h)\geq\mathcal N_N(h_{\mathrm{FS}})
\]
for every admissible \(h\) and each fixed \(N\), with the same normalization. A large-\(N\) asymptotic minimization is a different question; the original sentence does not distinguish them.

Under the aid's Gaussian-almost-sure simple-zero and Morse hypotheses, let \(M\) be the number of local maxima and \(S\) the number of saddles of the section norm away from its \(N\) zeros. It reports
\[
S-M=N-2,\qquad M\geq1,\qquad
\#\operatorname{Crit}=N-2+2M\geq N.
\]
Thus, in the stated fixed-degree setting, minimizing expected count is equivalent to minimizing expected maxima whenever these expectations are defined, as emphasized for \(N\geq2\) by the aid. The aid reports that the exact Fubini--Study count attains the lower bound for \(N=1,2\), giving global minimization in those degrees. It expressly does not assert uniqueness: constant metric scaling and \(\operatorname{PGL}(2,\C)\) pullbacks yield ties.

A separate asymptotic/local calculation in the aid concerns an exactly area-preserving conformal variation at Fubini--Study. With a real mean-zero variation \(f\), it records the constrained Calabi Hessian
\[
\frac12\int (Lf-8f)^2
\]
and the corresponding leading expected-count Hessian
\[
\frac1{27\pi N}\int (Lf-8f)^2.
\]
Here \(L\), the volume normalization, and the precise curvature-square functional must use the aid's conventions, which are not fully stated in the record; it identifies the kernel on area-preserving variations as the degree-one spherical harmonics. This leading large-\(N\) expression is not an exact finite-degree Hessian formula or a proof of global minimization for every \(N\). The aid describes exact fixed-degree global minimization for \(N\geq3\) as unresolved in the literature it checked. All supplied formulas and status claims remain unreviewed research-aid content.
\subsection{Short English statement}
Does Fubini--Study minimize an expected critical-point count on \(\PP^1\)? The source omits the random model. Under the supplied Gaussian holomorphic-section reconstruction, the aid reports a Morse-theoretic lower bound attained in degrees one and two, while distinguishing local asymptotic Hessians from the unresolved higher-degree global question.
\subsection{Sources}
\begin{itemize}
\item Bernard Shiffman, Question 6, AIM workshop \emph{Random analytic functions}: \url{https://aimath.org/WWN/randomzeros/arcc1.pdf}.
\item The supplied verification attributes the reconstructed ensemble to papers of Douglas--Shiffman--Zelditch, but provides no titles, years, or complete bibliography. It explicitly says the PDF itself supplies no degree, law, or critical-point definition.
\item Local registry \texttt{data/unsolvedmath/problems.json}, permanent ID \(20000631\), exact source code \texttt{catalogue entry}; canonical \texttt{aim-analysis-notes.json}, notes[162].
\item The attached research attempt supplies the Morse reduction, special degrees, and Hessian claims. It is explicitly machine-generated, not peer reviewed, and requires independent verification before citation.
\end{itemize}
\end{document}
